\documentclass{article}
\usepackage{comment}
\usepackage{amsmath}
\usepackage{amsfonts}
\usepackage{sectsty}
\usepackage{hyperref}
\usepackage{fancyhdr}
\usepackage[top=1in,bottom=1in,left=1in,right=1in]{geometry}
\usepackage{float}
\usepackage{graphicx}
\usepackage{tabularx}
\usepackage{mathtools}

\title{Assignment 1}
\author{Matthew Farah, \href{mailto:farahm11@mcmaster.ca}{$\langle$farahm11@mcmaster.ca$\rangle$}, 400468251}
\date{\today}

\hypersetup{
        colorlinks=true,
        linkcolor=blue,
        filecolor=magenta,
        urlcolor=blue,
        citecolor=blue,
	anchorcolor=blue
}

\fancyhead{}
\fancyhead[L]{Matthew Farah, \href{mailto:farahm11@mcmaster.ca}{$\langle$farahm11@mcmaster.ca$\rangle$}, 400468251}
\fancyhead[R]{Assignment 1}

\newenvironment{directsubsubs}{
  \makeatletter
  \renewcommand\thesubsubsection{\thesection.\Roman{subsubsection}}
  \makeatother
}{
  \makeatletter
  \renewcommand\thesubsubsection{\thesubsection.\Roman{subsubsection}}
  \makeatother
}

\renewcommand{\thesubsection}{\thesection.\alph{subsection}}
\renewcommand{\thesubsubsection}{\thesection.\alph{subsection}.\Roman{subsubsection}}
\makeatletter
\DeclareFontFamily{U}{MnSymbolA}{}
\DeclareFontShape{U}{MnSymbolA}{m}{n}{
    <-6>  MnSymbolA5
   <6-7>  MnSymbolA6
   <7-8>  MnSymbolA7
   <8-9>  MnSymbolA8
   <9-10> MnSymbolA9
  <10-12> MnSymbolA10
  <12->   MnSymbolA12}{}
\DeclareFontShape{U}{MnSymbolA}{b}{n}{
    <-6>  MnSymbolA-Bold5
   <6-7>  MnSymbolA-Bold6
   <7-8>  MnSymbolA-Bold7
<8-9>  MnSymbolA-Bold8
   <9-10> MnSymbolA-Bold9
  <10-12> MnSymbolA-Bold10
  <12->   MnSymbolA-Bold12}{}
\DeclareSymbolFont{MnSyA}{U}{MnSymbolA}{m}{n}
\SetSymbolFont{MnSyA}{bold}{U}{MnSymbolA}{b}{n}

\DeclareRobustCommand{\overleftharpoon}{\mathpalette{\overarrow@\leftharpoonfill@}}
\DeclareRobustCommand{\overrightharpoon}{\mathpalette{\overarrow@\rightharpoonfill@}}
\def\leftharpoonfill@{\arrowfill@\leftharpoondown\mn@relbar\mn@relbar}
\def\rightharpoonfill@{\arrowfill@\mn@relbar\mn@relbar\rightharpoonup}

\DeclareMathSymbol{\leftharpoondown}{\mathrel}{MnSyA}{'112}
\DeclareMathSymbol{\rightharpoonup}{\mathrel}{MnSyA}{'100}
\DeclareMathSymbol{\mn@relbar}{\mathrel}{MnSyA}{'320}
\makeatother


\newcolumntype{Y}{>{\centering\arraybackslash}X}
\renewcommand{\arraystretch}{1.8}

\sectionfont{\fontsize{12}{12}\bfseries}
\subsectionfont{\fontsize{10}{12}\normalfont}
\subsubsectionfont{\fontsize{10}{12}\normalfont}

\newcommand{\harpoon}[1]{\overrightharpoonup{#1}}

\begin{document}
\maketitle
\pagestyle{fancy}

\begin{center}
	\Large{\underline{Discrete Systems}}
\end{center}
\section[Question ~\thesection]{The system S is given by $y(n) = y(n - 1) - x(n - 1) + x(n)$. Assume that $y(-1) = 0$}

\subsection[~\thesubsection]{Compute the output to the input $x(n) = \delta(n)$}

\begin{tabularx}{\linewidth}{| c || Y Y Y Y Y Y |}
	\hline
	$n$ & 0 & 1 & 2 & 3 & 4 & 5 \\
	\hline
	$x(n)$ & 1 & 0 & 0 & 0 & 0 & 0 \\
	$y(n)$ & 1 & 0 & 0 & 0 & 0 & 0 \\
	\hline
\end{tabularx}

\subsection[~\thesubsection]{Compute the output to the input $x(n) = \delta(n) + \delta(n - 1) + \delta(n - 2)$}

\begin{tabularx}{\linewidth}{| c || Y Y Y Y Y Y |}
	\hline
	$n$ & 0 & 1 & 2 & 3 & 4 & 5 \\
	\hline
	$x(n)$ & 1 & 1 & 1 & 0 & 0 & 0 \\
	$y(n)$ & 1 & 1 & 1 & 0 & 0 & 0 \\
	\hline
\end{tabularx}

\subsection[~\thesubsection]{Compute the output to the input $x(n) = \cos({\pi}n)$}

\begin{tabularx}{\linewidth}{| c || Y Y Y Y Y Y |}
	\hline
	$n$ & 0 & 1 & 2 & 3 & 4 & 5 \\
	\hline
	$x(n)$ & 1 & -1 & 1 & -1 & 1 & -1 \\
	$y(n)$ & 2 & 0 & 2 & 0 & 2 & 0 \\
	\hline
\end{tabularx}

\newpage

\begin{center}
	\Large{\underline{Difference Equations to State Space Equations}}
\end{center}

\section[Question ~\thesection]{Given the difference equation $y(n) = y(n - 1) - y(n - 2) + x(n) - x(n - 2)$, determine the $(A, B, C, D)$ representation of the system.}

\begin{directsubsubs}
	
	\subsubsection[~\thesubsubsection]{Find $\overrightharpoon{S(n)}$}
	
	\begin{align*}
		&\text{Let:}& \text{Then:} \\
		&s_1(n) = x(n - 2) &s_1(n + 1) &= x(n - 1) = s_2(n) \\
		&s_2(n) = x(n - 1) &s_2(n + 1) &= x(n) \\
		&s_3(n) = y(n - 2) &s_3(n + 1) &= y(n - 1) = s_4(n) \\
		&s_4(n) = y(n - 1) &s_4(n + 1) &= y(n) \\
		& & &= y(n - 1) - y(n - 2) + x(n) - x(n - 2) \\
		& & &= -s_1(n) - s_3(n) + s_4(n) + x(n)
	\end{align*}

	Thus, let $\overrightharpoon{S(n)} = \begin{psmallmatrix} s_1 \\ s_2 \\ s_3 \\ s_4 \end{psmallmatrix}$

	\subsubsection[~\thesubsubsection]{Find $\overrightharpoon{S(n + 1)}$}
	
	\begin{align*}
		\overrightharpoon{S(n + 1)} &= \begin{pmatrix} s_1(n + 1) \\ s_2(n + 1) \\ s_3(n + 1) \\ s_4(n + 1) \end{pmatrix} \\
		&= \begin{pmatrix} s_2(n) \\ x(n) \\ s_4(n) \\ -s_1(n) - s_3(n) + s_4(n) + x(n) \end{pmatrix} \\
		&= \begin{pmatrix} 0 & 1 & 0 & 0 \\ 0 & 0 & 0 & 0 \\ 0 & 0 & 0 & 1 \\ -1 & 0 & -1 & 1 \end{pmatrix}\overrightharpoon{S(n)} + \begin{pmatrix} 0 \\ 1 \\ 0 \\ 1 \end{pmatrix}x(n)
	\end{align*}

	Therefore, our $\boldsymbol{A}$ and $\boldsymbol{B}$ matrices are $\begin{psmallmatrix} 0 & 1 & 0 & 0 \\ 0 & 0 & 0 & 0 \\ 0 & 0 & 0 & 1 \\ -1 & 0 & -1 & 1 \end{psmallmatrix}$ and $\begin{psmallmatrix} 0 \\ 1 \\ 0 \\ 1 \end{psmallmatrix}$.

	\subsubsection[~\thesubsubsection]{Find $y(n)$}

	\begin{align*}
		y(n) &= -x(n - 2) - y(n - 2) + y(n + 1) + x(n) \\
		&= -s_1(n) -s_3(n) + s_4(n) + x(n) \\
		&= \begin{pmatrix} -1 & 0 & -1 & 1 \end{pmatrix}\overrightharpoon{S(n)} + \begin{pmatrix} 1 \end{pmatrix}x(n)
	\end{align*}

	Therefore, out $\boldsymbol{C}$ and $\boldsymbol{D}$ matrices are $\begin{psmallmatrix} -1 & 0 & -1 & 1 \end{psmallmatrix}$ and $\begin{psmallmatrix} 1 \end{psmallmatrix}$ respectively.

	Putting it all together, our system can be represented in state space form as:
	
	\begin{align*}
		\overrightharpoon{S(n + 1)} &= \begin{pmatrix} 0 & 1 & 0 & 0 \\ 0 & 0 & 0 & 0 \\ 0 & 0 & 0 & 1 \\ -1 & 0 & -1 & 1 \end{pmatrix}\overrightharpoon{S(n )} + \begin{pmatrix} 0 \\ 1 \\ 0 \\ 1 \end{pmatrix}x(n) \\\\
		y(n) &= \begin{pmatrix} -1 & 0 & -1 & 1 \end{pmatrix}\overrightharpoon{S(n)} + \begin{pmatrix} 1 \end{pmatrix}x(n)
	\end{align*}

\end{directsubsubs}

\newpage

\begin{center}
	\Large{\underline{$(A, B, C, D)$ Representation}}
\end{center}

\section[Question ~\thesection]{Given the system $y(n) = x(n) + x(n - 2) - y(n - 2)$}

\subsection[~\thesubsection]{Find the $(A, B, C, D)$ representation of the system}

\subsubsection[~\thesubsubsection]{Find $\overrightharpoon{S(n)}$}

\begin{align*}
       &\text{Let:}& \text{Then:} \\
       &s_1(n) = x(n - 2) &s_1(n + 1) &= x(n - 1) = s_2(n) \\
       &s_2(n) = x(n - 1) &s_2(n + 1) &= x(n) \\
       &s_3(n) = y(n - 2) &s_3(n + 1) &= y(n - 1) = s_4(n) \\
       &s_4(n) = y(n - 1) &s_4(n + 1) &= y(n) \\
       & & &= -y(n - 2) + x(n - 2) + x(n) \\
       & & &= s_1(n) - s_3(n) + x(n)
\end{align*}

Thus, let $\overrightharpoon{S(n)} = \begin{psmallmatrix} s_1(n) \\ s_2(n) \\ s_3(n) \\ s_4(n) \end{psmallmatrix}$.

\subsubsection[~\thesubsubsection]{Find $\overrightharpoon{S(n + 1)}$}

\begin{align*}
	\overrightharpoon{S(n+1)} &= \begin{pmatrix} s_1(n + 1) \\ s_2(n + 1) \\ s_3(n + 1) \\ s_4(n + 1) \end{pmatrix} \\
	&= \begin{pmatrix} s_2(n) \\ x(n) \\ s_4(n) \\ s_1(n) - s_3(n) + x(n) \end{pmatrix} \\
	&= \begin{pmatrix} 0 & 1 & 0 & 0 \\ 0 & 0 & 0 & 0 \\ 0 & 0 & 0 & 1 \\ 1 & 0 & -1 & 0 \end{pmatrix}\overrightharpoon{S(n)} + \begin{pmatrix} 0 \\ 1 \\ 0 \\ 1 \end{pmatrix}
\end{align*}

Therefore, our $\boldsymbol{A}$ and $\boldsymbol{B}$ matrices are $\begin{psmallmatrix} 0 & 1 & 0 & 0 \\ 0 & 0 & 0 & 0 \\ 0 & 0 & 0 & 1 \\ 1 & 0 & -1 & 0 \end{psmallmatrix}$ and $\begin{psmallmatrix} 0 \\ 1 \\ 0 \\ 1 \end{psmallmatrix}$ respectively.

\subsubsection[~\thesubsubsection]{Find $y(n)$}

\begin{align*}
	y(n) &= -y(n - 2) + x(n - 2) + x(n) \\
	&= s_1(n) - s_3(n) + x(n) \\
	&= \begin{pmatrix} 1 & 0 & -1 & 0 \end{pmatrix}\overrightharpoon{S(n)} + \begin{pmatrix} 1 \end{pmatrix}x(n)
\end{align*}

Therefore, our $\boldsymbol{C}$ and $\boldsymbol{D}$ matrices are $\begin{psmallmatrix} 1 & 0 & -1 & 0 \end{psmallmatrix}$ and $\begin{psmallmatrix} 1 \end{psmallmatrix}$ respectively.

Putting it all together, our system can be represented in state space form as:

\begin{align*}
	\overrightharpoon{S(n + 1)} &= \begin{pmatrix} 0 & 1 & 0 & 0 \\ 0 & 0 & 0 & 0 \\ 0 & 0 & 0 & 1 \\ 1 & 0 & -1 & 0 \end{pmatrix}\overrightharpoon{S(n)} + \begin{pmatrix} 0 \\ 1 \\ 0 \\ 1 \end{pmatrix} \\
	y(n) &= \begin{pmatrix} 1 & 0 & -1 & 0 \end{pmatrix}\overrightharpoon{S(n)} + \begin{pmatrix} 1 \end{pmatrix}x(n)
\end{align*}

\subsection[~\thesubsection]{Use the state space equations to compute 5 outputs to the input $x(n) = \delta(n)$}

\begin{align*}
	&\underline{\text{Current State} \; (\overrightharpoon{S(n)})}& &\underline{\text{Current Output} \; (y(n))} \\ \\
	&\overrightharpoon{S(0)} = \overrightharpoon{0}& &y(0) = \begin{pmatrix} 1 & 0 & -1 & 0 \end{pmatrix}\overrightharpoon{S(0)} + \begin{pmatrix} 1 \end{pmatrix}x(0) = 1 \\
	&\overrightharpoon{S(1)} = \begin{pmatrix} 0 & 1 & 0 & 0 \\ 0 & 0 & 0 & 0 \\ 0 & 0 & 0 & 1 \\ 1 & 0 & -1 & 0 \end{pmatrix}\overrightharpoon{S(0)} + \begin{pmatrix} 0 \\ 1 \\ 0 \\ 1 \end{pmatrix}x(0) = \begin{pmatrix} 0 \\ 1 \\ 0 \\ 1 \end{pmatrix}& &y(1) = \begin{pmatrix} 1 & 0 & -1 & 0 \end{pmatrix}\overrightharpoon{S(1)} + \begin{pmatrix} 1 \end{pmatrix}x(1) = 0 \\
	&\overrightharpoon{S(2)} = \begin{pmatrix} 0 & 1 & 0 & 0 \\ 0 & 0 & 0 & 0 \\ 0 & 0 & 0 & 1 \\ 1 & 0 & -1 & 0 \end{pmatrix}\overrightharpoon{S(1)} + \begin{pmatrix} 0 \\ 1 \\ 0 \\ 1 \end{pmatrix}x(1) = \overrightharpoon{0}& &y(2) = \begin{pmatrix} 1 & 0 & -1 & 0 \end{pmatrix}\overrightharpoon{S(2)} + \begin{pmatrix} 1 \end{pmatrix}x(2) = 0 \\
	&\overrightharpoon{S(3)} = \begin{pmatrix} 0 & 1 & 0 & 0 \\ 0 & 0 & 0 & 0 \\ 0 & 0 & 0 & 1 \\ 1 & 0 & -1 & 0 \end{pmatrix}\overrightharpoon{S(2)} + \begin{pmatrix} 0 \\ 1 \\ 0 \\ 1 \end{pmatrix}x(2) = \overrightharpoon{0}& &y(3) = \begin{pmatrix} 1 & 0 & -1 & 0 \end{pmatrix}\overrightharpoon{S(3)} + \begin{pmatrix} 1 \end{pmatrix}x(3) = 0 \\
	&\overrightharpoon{S(4)} = \begin{pmatrix} 0 & 1 & 0 & 0 \\ 0 & 0 & 0 & 0 \\ 0 & 0 & 0 & 1 \\ 1 & 0 & -1 & 0 \end{pmatrix}\overrightharpoon{S(3)} + \begin{pmatrix} 0 \\ 1 \\ 0 \\ 1 \end{pmatrix}x(3) = \overrightharpoon{0}& &y(4) = \begin{pmatrix} 1 & 0 & -1 & 0 \end{pmatrix}\overrightharpoon{S(4)} + \begin{pmatrix} 1 \end{pmatrix}x(4) = 0 \\
\end{align*}

\newpage

\begin{center}
	\large{\underline{Compound Interest}}
\end{center}

\section[Question ~\thesection]{A bank account is just a system: Letting $x(n)$ represent the amount deposited or withdrawn (with $x(n) > 0$ representing deposits and $x(n) < 0$ representing withdrawals), and $\alpha$ represent the interest rate, we can express the account balance with the system $y(n) = y(n - 1) + {\alpha}y(n - 1) + x(n)$. How much money do you have after 10 days if you initially deposit \$100?}

We're just going to solve this through direct computation with a relatively simply inference about $y(n)$:

\vspace{1cm}

\begin{tabularx}{\linewidth}{| c || Y Y Y Y Y Y |}
	\hline
	$n$ & 0 & 1 & 2 & 3 & $\cdots$ & 10 \\
	\hline
	$x(n)$ & 100 & 0 & 0 & 0 & $\cdots$ & 0 \\
	$y(n)$ & 100 & $100(1 + \alpha)$ & $100(1 + \alpha)^2$ & $100(1 + \alpha)^3$ & $\cdots$ & $100(1 + \alpha)^{10}$ \\
	\hline
\end{tabularx}

\vspace{1cm}

Therefore, with an initial deposit of \$100, the account balance grows to \$$100(1 + \alpha)^{10}$ after 10 days.

\end{document}
